\DTMsavedate{mydate}{2022-06-30}
\maketitle{Engineering Mathematics \Romannum{2}}{Applied Mathematics}{\DTMusedate{mydate}}

% ----------------------------------------------------- %
% COURSE INFO
\subsection*{Course Information}\label{course:sma2116}
\vspace{0ex}%
\noindent\parbox{0.5\textwidth}{%
    \noindent\begin{tabular}{@{}ll}
                 \textsf{Course:}          & 	Engineering Mathematics \Romannum{2}\\
                 \textsf{Course Code:}         & 	SMA 2116    \\
                 \textsf{Credit Hours:}    & 	10
    \end{tabular}}
\vspace{2ex}\hrule\vspace{2ex}

% ----------------------------------------------------- %
% INSTRUCTOR INFO
\subsection*{Lecturer Information}
\noindent\parbox{0.5\textwidth}{%
    \noindent\begin{tabular}{@{}ll}
                 \textsf{Name:}         &    Mlamuli Dhlamini           \\
% \textsf{Phone:}        &    \phone          \\
\textsf{Email:}        &    mlamuli.dhlamini@nust.ac.zw    \\
\textsf{Qualifications:}        &  PhD in Applied Mathematics, MSc in Industrial Mathematics, BSc in Applied \\
 & Mathematics
    \end{tabular}}
\vspace{2ex}\hrule\vspace{2ex}


% ----------------------------------------------------- %
% COURSE DESCRIPTION
\subsection*{Learning Aim}
This course aims at introducing students to dealing with functions of more than one variable and Fourier series.
\vspace{2ex} \hrule\vspace{2ex}

% ----------------------------------------------------- %
% LEARNING OBJECTIVES
\subsection*{Learning Objectives}
At the end of the course students must be able to:
\begin{outline}
    \1 evaluate multiple integrals.

    \1 change from one co-ordinate system to another.

    \1 use the integral theorems.

    \1 to evaluate the fourier series for a given function.

\end{outline}
\vspace{2ex}\hrule\vspace{2ex}

\subsection*{Learning Requirements}
Students are expected to have completed and passed the following courses:
\begin{outline}
    \1      SMA 1116

    \1      SMA 1216

\end{outline}
\vspace{2ex}\hrule\vspace{2ex}

% ----------------------------------------------------- %
% COURSE TOPICS
\subsection*{Topics}
\begin{outline}
    \1 Multiple Integrals
    \1 Vector Calculus
    \1 Vector Integral Function
    \1 Integral Theorem
    \1 Fourier’s Series
\end{outline}
\vspace{2ex} \hrule\vspace{2ex}

% ----------------------------------------------------- %
% Students Assenssment
\subsection*{Assessment of Students}
At least two in-class tests will be given as a continuous assessment during the semester before a final exam at the end of the semester.
\vspace{2ex}\hrule\vspace{2ex}

% ----------------------------------------------------- %
% Grade Evaluation
\subsection*{Course Evaluation and Grading Policies}
Grades are based on:
Continuous assessment tests (\qty{25}{\percent}),
Final Exam (\qty{75}{\percent})
\vspace{2ex}\hrule\vspace{2ex}
% ----------------------------------------------------- %
% EQUIPMENT
% https://sites.udel.edu/computing-purchases/personal-specs/
% \subsection*{Essential Equipment}

% \vspace{2ex}\hrule\vspace{2ex}

% Policies and Other information
\subsection*{Policies and Other Information}
% Key Text or some Books
\subsection*{Key Text}
\begin{itemize}
    \item Calculus with analytic geometry by Dennis Zill
    \item Vector Analysis Schaum’s outline Series by Murray R Spiegel
    \item Vector calculus by Miroslav Lovric
    \item Elementary applied Partial differential equations by Richard Haberman
    \item Introduction to Fourier Analysis by Norman Morrison
\end{itemize}
\vspace{2ex}\hrule
\vspace{2ex}\hrule\vspace{2ex}